\documentclass[UTF8,12pt]{ctexart}
\usepackage[a4paper,margin=24mm]{geometry}
\usepackage{amsmath,amssymb,bm,graphicx,adjustbox}
\newcommand{\fitmath}[1]{\begin{adjustbox}{max width=\linewidth}$\displaystyle #1$\end{adjustbox}}
\setlength{\parindent}{0pt}
\setlength{\parskip}{.7em}
\sloppy
\allowdisplaybreaks
\begin{document}
\section*{2013 年考研数学三 · 参考答案}
\subsection*{原 PDF 第 35 页}

\subsection*{2013 年真题参考答案}

\subsection*{一、选择题}

(1) D。 (2) C。 (3) B。 (4) D。 (5) B。 (6) B。 (7) A。 (8) C。


\subsection*{二、填空题}

(9) \(\displaystyle -2\)。 (10) \(\displaystyle 2(1-\ln2)\)。 (11) \(\displaystyle \ln2\)。 (12) \(\displaystyle (C_1+C_2x)e^{x/2}\)。 (13) \(\displaystyle -1\)。 (14) \(\displaystyle 2e^2\)。


\subsection*{三、解答题}

(15) \(\displaystyle n=2,\allowbreak \ a=7\)。


(16) \(\displaystyle a=7\sqrt7\)。


(17) \(\displaystyle \dfrac{\displaystyle 416}{\displaystyle 3}\)。


(18)（I）\(\displaystyle L^{\prime}(Q)=40-\dfrac{\displaystyle Q}{\displaystyle 500}\)。（II）当 \(\displaystyle p=50\) 时，边际利润为 20。其经济意义为：当 \(\displaystyle p=50\) 时，销售第 10 001 件商品时所得的利润为 20 元。（III）当定价为 40 元时，利润最大。


(19) 证明略。


(20) \(\displaystyle a=-1,\allowbreak b=0\) 时，\(\displaystyle C=\begin{pmatrix}k_1+k_2+1&-k_2\\k_2&k_1\end{pmatrix}\)，其中 \(\displaystyle k_1,\allowbreak k_2\) 为任意常数。


(21) 证明略。


(22)（I）\(\displaystyle f(x,\allowbreak y)=\begin{cases}\dfrac{\displaystyle 9y^2}{\displaystyle x},\allowbreak &0<y<x,\allowbreak \ 0<x<1,\allowbreak \\0,\allowbreak &\text{其他}\end{cases}\)。（II）\(\displaystyle f_Y(y)=\begin{cases}-9y^2\ln y,\allowbreak &0<y<1,\allowbreak \\0,\allowbreak &\text{其他}\end{cases}\)。（III）\(\displaystyle \dfrac{\displaystyle 1}{\displaystyle 8}\)。


(23)（I）\(\displaystyle \hat\theta=\overline X\)。（II）\(\displaystyle \hat\theta=\dfrac{\displaystyle 2n}{\displaystyle \sum_{i=1}^n1/X_i}\)。


\clearpage
\end{document}
